
%\documentclass[oneside,final,14pt]{extreport}
\documentclass[12pt, a4paper, openany]{book}



\usepackage[left=2cm,right=2cm,top=2cm,bottom=2cm,bindingoffset=0cm]{geometry}
%\usepackage[koi8-r]{inputenc}
%\usepackage[russianb]{babel}
\usepackage{vmargin}

\setpapersize{A4}
\usepackage[T2A]{fontenc}
\usepackage[utf8x]{inputenc}
\usepackage[english, russian]{babel}
\setmarginsrb{2cm}{2cm}{2cm}{2cm}{0pt}{5mm}{0pt}{0mm}
\usepackage{indentfirst}
\usepackage{nicefrac} % For comparison
%\usepackage{xfrac}    % Works better with other fonts
%\usepackage[unicode, pdftex]{hyperref}
\usepackage{lettrine}
\usepackage[usenames]{color}
\usepackage{colortbl}
\usepackage{mathtext}
\usepackage{epigraph}
\usepackage{amsmath, amsfonts, amssymb, mathrsfs}
%\usepackage{mathptmx}
%\usepackage{txfonts}
\usepackage{pxfonts}
\usepackage[pagestyles]{titlesec}
\usepackage{ebgaramond}
\usepackage{awesomebox}
\usepackage{enumitem}
\usepackage{makeidx}
\usepackage[letterspace=150]{microtype}
\makeindex
\usepackage{etoolbox}
\makeatletter
\newlength\epitextskip
\pretocmd{\@epitext}{\em}{}{}
\apptocmd{\@epitext}{\em}{}{}
\patchcmd{\epigraph}{\@epitext{#1}\\}{\@epitext{#1}\\[\epitextskip]}{}{}
\makeatother
%running fraction with slash - requires math mode.
\newcommand*\rfrac[2]{{}^{#1}\!/_{#2}}

\DeclareSymbolFont{Xlargesymbols}{OMX}{cmex}{m}{n}
\DeclareMathSymbol{\Xsum}{\mathop}{Xlargesymbols}{80}

\setlength\epigraphrule{0pt}
\setlength\epitextskip{2ex}
\setlength\epigraphwidth{.8\textwidth}

\usepackage{xfrac}    % Works better with other fonts
\usepackage[colorlinks=true,linkcolor=black,urlcolor=black,bookmarksopen=true]{hyperref}

\usepackage{fancyhdr} % пакет для установки колонтитулов
\pagestyle{fancy} % смена стиля оформления страниц
\fancyhf{} % очистка текущих значений
\fancyhead[C]{\thepage} % установка верхнего колонтитула
\renewcommand{\headrulewidth}{0pt} % убрать разделительную линию


% Настройка вертикальных и горизонтальных отступов
\titlespacing{\chapter}{0pt}{5pt}{5pt}
\titlespacing{\section}{\parindent}{4mm}{4mm}
\titlespacing{\subsection}{\parindent}{3mm}{3mm}


%
%\renewcommand{\tabcolsep}{1cm}   %% increase table column spacing

% Настройка задачи со зведочкой
\newcounter{namedlistcounter}  % number the items
\newenvironment{withdot}
{\begin{list}
		{\arabic{namedlistcounter}*.} % labeling 
		{\usecounter{namedlistcounter}   % set counter
			\setlength{\leftmargin}{3em}} % set spacing 
	}
	{\end{list}}


\newcommand{\anonsection}[1]{ \section*{#1} \addcontentsline{toc}{section}{\numberline {}#1}} 

\makeatletter %%%%% <---- Starting chapter without a pagebreak
\renewcommand\chapter{\par%
	\thispagestyle{plain}% \global\@topnum\z@
	\@afterindentfalse \secdef\@chapter\@schapter}
\makeatother %%%%% <---- Starting chapter without a pagebreak
\titleformat{\chapter}[display]
{\normalfont\bfseries}{}{0pt}{\Large}

\newpagestyle{mystyle}{
	\sethead[\thepage][][]{}{}{\thepage}	
}

\renewcommand{\rmdefault}{cmr}


\pagestyle{mystyle}
\sloppy
\begin{document}


\begin{titlepage}

\begin{center}
\topskip0pt
\vspace*{\fill}


{\large\bf В. Л. ПЯНКЕВИЧ\\}
\ \\
\ \\
{\Huge\bf ЛЮДИ ЖИЛИ СЛУХАМИ\\}
\ \\
{\Large \bf НЕФОРМАЛЬНОЕ КОММУНИКАТИВНОЕ ПРОСТРАНСТВО БЛОКАДНОГО ЛЕНИНГРАДА}
\ \\
\ \\
\vspace*{\fill}

\vfill


Санкт-Петербург

«Владимир Даль»

2014


\end{center}

\end{titlepage}

\thispagestyle{empty} % выключаем отображение номера для этой страницы

\newpage


\setcounter{secnumdepth}{0}

\epigraph{Конечно, слухи. Я, как всегда, не записываю всего вздора, который болтают. Это мусор, сор, слякоть обывательщины. Но иногда среди всяких сплетен прорывается какая-нибудь мысль, предположение, которое - как дым не без oгня.}{\textit{Г. А. Князев}}

\section[Введение]{\center \textbf{ВВЕДЕНИЕ}}

Накануне и во время блокады Ленинграда неформальные коммуникации занимали очень важное место в жизни горожан. Сознание сотен тысяч людей, существовавших в условиях постоянной угрозы гибели от голода, холода, немецкой бомбы или снаряда, находило выражение в самых разных формах стихийного общения: рассказах очевидцев, слухах, письмах, вопросах и обращениях к власти, спонтанных высказываниях, частных объявлениях о продаже и покупке продуктов питания, вещей, с предложением услуг и многом другом. Такие формы автокоммуникации, как мечты и сны, рефлексии, внутренние монологи, фиксировались в дневниковых записях ленинградцев. Дневники, по сути, исповеди, в большинстве своем не предназначались для чужих глаз, хотя порой автор понимал и даже надеялся, что его печальные записи увидят потомки.

Люди говорили, писали, думали о смертельной нехватке продовольствия и карточных нормах, угрозе вторжения немцев в город и прорыве блокады, о самих себе, согражданах, врагах, власти, прошлом и будущем. Все это в свою очередь определяло эмоциональное состояние и поведение сотен тысяч людей в катастрофических условиях блокады. Молва предваряла наступление события, сопровождала его и сохраняла его образ в памяти людей, представляя собой устное народное творчество, городской фольклор. Стихийное информационное пространство было той призмой, через которую блокадники воспринимали реальность.

По наблюдению выдающегося физиолога И. П. Павлова, слово для русскою человека всегда имело приоритетное значение:



\newpage
\tableofcontents



\thispagestyle{empty} %


\newpage

\setcounter{secnumdepth}{0}



\phantomsection
\begin{center}
	\ \\
	\ \\
	\ \\
	\ \\
	\ \\
	\ \\

	\ \\
	\ \\
	\ \\
	\ \\
	\ \\
	\ \\
	\ \\
	\ \\
	\ \\
\textit{Научное издание}
	\ \\
	\ \\
	\ \\
	\ \\
	\textbf{Владимир Леонидович Пянкевич}
	\ \\
	\ \\
«\textbf{ЛЮДИ ЖИЛИ СЛУХАМИ}»:

Неформальное коммуникативное пространство блокадного Ленинграда
	\ \\
	\ \\
	\ \\		
	Редактор издательство \textit{Т. В. Глушенкова}
	
	Редактор \textit{Т. Д. Рахманова}
	
	Компьютерная верстка \textit{Л. А. Шитовой}
	
	\ \\
	\ \\
	
	\ \\	

	\ \\
\ \\

\ \\
	\ \\	

\ \\
\ \\

\ \\		
	
		\ \\	
	
{\footnotesize 	 
	
	Формат 60$\times$184$\rfrac{1}{16}$. Гарнитура Таймс.
	
	Подписано к печати 20.09.2013.
	
	Бумага офсетная. Печать офсетная. Усл. печ. л. 30,0.
	
	Уч.-изд. л. 26,7 + вкл. Тип. зак. № 3027
		\ \\
				\ \\
		
Издательство «Владимир Даль»

193036, Санкт-Петербург, ул. 7-я Советская, д. 19
	
	\ \\
	
Первая Академическая типография «Наука»

199034, Санкт-Петербург, 9 линия, 12}


\end{center}
\thispagestyle{empty} % выключаем отображение номера для этой страницы



\newpage

\setcounter{secnumdepth}{0}

\phantomsection

\section*{Описание}

{\bf Название:} Крупнейшие уголовные дела ХХ века в США

{\bf Автор:} Валентин Алексеевич Ковалев

{\bf Издательство:} М.: Юрид. лит., 1990

{\bf ISNB} 5-7260-0260-1

{\bf Аннотация:} Эта книга — документальное повествование о крупнейших в  современной истории США судебных процессах по обвинению  американских граждан в совершении общеуголовных и уголовно-политических преступлений. Она знакомит читателя с подлинными материалами уголовных дел о вооруженных нападениях и убийствах, организации взрывов и похищении людей, антиправительственных заговорах и  атомном шпионаже. 

Вокруг этих дел бушевали политические страсти и юридические дискуссии. Наряду с профессиональными юристами в них принимали участие видные политические и государственные деятели: ученые, сенаторы, министры, кандидаты в президенты, президенты США. 

Для широкого круга читателей. 
\thispagestyle{empty} % выключаем отображение номера для этой страницы


\newpage

\setcounter{secnumdepth}{0}  

\phantomsection	

\ \\
\ \\
\ \\
\ \\

\hangindent=2cm \hangafter=0  	\textbf{ИНФОРМАЦИЯ ОБ ОЦИФРОВКЕ}

\hangindent=2cm \hangafter=0  	\textbf{Источник}:  \href{https://rutracker.org/forum/viewtopic.php?t=3985675}{https://rutracker.org/forum/viewtopic.php?t=3985675} 

\hangindent=2cm \hangafter=0  	\textbf{Версия}: 1.0 - от 15 июля 2026

\hangindent=2cm \hangafter=0  	\textbf{Сканирование}: Skaramusch

\hangindent=2cm \hangafter=0  	\textbf{OCR, DjVuing}: ivanstor (2012)

\hangindent=2cm \hangafter=0  	\textbf{Оцифровка}: lord199 (\href{mailto:lord199@mail.ru}{lord199@mail.ru})



\hangindent=2cm \hangafter=0  Если вы нашли какие-то опечатки и другие неточности, просьба связаться по электронному адресу.




\thispagestyle{empty} % выключаем отображение номера для этой страницы


\end{document}